\documentclass{easychair}

\usepackage{amsmath,amssymb,amsthm}

\newtheorem{definition}{Definition}
\newtheorem{theorem}{Theorem}
\newtheorem{proposition}{Proposition}

\title{Fuzzifying Derivatives of Regular Expressions}

\author{
Mariam Katamashvili\inst{1,2}
\and David Maisuradze\inst{1}
}

\institute{
Kutaisi International University, Kutaisi, Georgia
\and
VIAM, Tbilisi State University, Tbilisi, Georgia
}
\authorrunning{Mariam Katamashvili, David Maisuradze}
\titlerunning{Fuzzifying Derivatives of Regular Expressions}

\begin{document}

\maketitle
There is growing interest in the computer science and AI communities in moving from {\em exact}, {\em crisp}, {\em qualitative} settings to {\em metric}, {\em probabilistic}, {\em quantitative}~\cite{AlvesKesnerRamos2025,SilvaBonchiBonsangueRutten2011} and hence ultimately {\em fuzzy} settings~\cite{met05}. In this view, we explore how to extend the standard theory of finite automata~\cite{C71}, based on exact equality of symbols, to a setting where symbol matching is fuzzy~\cite{DHMM}. Namely, following similarity-based reasoning~\cite{Sessa2002},
we assume that the alphabet is enriched with a {\em similarity relation}
and that matching is performed up to {\em cut} values. In this extended abstract, we generalize the theory of
Brzozowski~\cite{B64} derivatives, further developed and revisited
in~\cite{OwensReppyTuron}, which gives a direct construction of automata
from regular expressions: after reading a word, the current state is represented by the
corresponding derivative of the expression. Most of our results are formalized in Rocq~\cite{Rocq}.

Let $\Sigma$ be a finite alphabet, and let
$\mathcal R:\Sigma\times\Sigma\to[0,1]$
be a {\em fuzzy similarity relation}, where $[0,1]$ is endowed with a
left-continuous T-norm, making it a residuated lattice~\cite{met05}.
Fix a cut value $0<\mu\le1$. A similarity relation satisfies reflexivity, symmetry, and the reverse triangle inequality with respect to the T-norm. In this extended abstract, we use G\"odel's minimum T-norm.

Similarity is extended from symbols to regular expressions, words, and languages.
\begin{definition}\hfill\\
$\bullet$ \textbf{Similarity on regular expressions and on words.}
For regular expressions $r_1,r_2\in E^\Sigma$ define
$
\mathcal R^e:E^\Sigma\times E^\Sigma\to[0,1]$
 by \[\mathcal R^{e}(r_1,r_2)=\begin{cases}\min_{1\le i\le n}\mathcal R(a_i,b_i) &\text{ if } r_1=r[a_1,\ldots,a_n]\text{ and }r_2=r[b_1,\ldots,b_n];
\\ 
0& \text{ otherwise. }
\end{cases}
\] where we denote $r[a_1,\ldots,a_n]$ to highlight the occurrences of the alphabet symbols $a_1,\ldots,a_n$ in $r$. Similarity on words $\mathcal R^w$ is derived from $\mathcal R^e$ viewing words as regular expressions.\\
$\bullet$ \textbf{Similarity on languages.}
For languages $L_1,L_2\subseteq\Sigma^\ast$, define
\[
\mathcal R^L(L_1,L_2)
=
\min\!\Bigl\{
\inf_{w_1\in L_1}\sup_{w_2\in L_2}\mathcal R^w(w_1,w_2),
\inf_{w_2\in L_2}\sup_{w_1\in L_1}\mathcal R^w(w_1,w_2)
\Bigr\}.
\]\end{definition}

The constructors $\&$ and $\neg$ in the definition of regular expressions, $E^\Sigma$, are not amenable to a well-behaved fuzzification {\em vis-à-vis} similarity on languages. Hence, the theorems are restricted to the fragment generated by $\emptyset,\varepsilon$,
letters, union, concatenation, and Kleene star, which we call {\em restricted regular expressions}.

\begin{proposition} The relations $\mathcal R^w$, $\mathcal R^L$, and $\mathcal R^e$ are similarity relations.
\end{proposition}

One of the main results of this work establishes a connection between our fuzzy analogue of the theory of finite automata and the crisp version of the theory for a corresponding quotient alphabet. Namely, we introduce the following definitions.
\begin{definition}For a fixed cut value $\mu$, define the following equivalence relations: \\
$\bullet$ $\boldsymbol{\sim_\mu^\Sigma}$ on the alphabet $\Sigma$ by

$a\sim_\mu^\Sigma b
\quad\Longleftrightarrow\quad
\mathcal R(a,b)\ge\mu;
$
\\
$\bullet$ $\boldsymbol{\sim_\mu^e}$ on regular expressions in $E^\Sigma$ by

$r\sim_\mu^e s
\quad\Longleftrightarrow\quad
\mathcal R^e(r,s)\ge\mu$;
equivalently, $r$ and $s$ are syntactically similar up to $\sim_\mu^\Sigma$-equivalent alphabet symbols.
\\
$\bullet$ $\boldsymbol{\sim_\mu^{\Sigma^\ast}}$ on words in $\Sigma^\ast$ by
 
$w_1\sim_\mu^{\Sigma^\ast} w_2
\quad\Longleftrightarrow\quad
\mathcal R^w(w_1,w_2)\ge\mu;
$\\
$\bullet$ $\boldsymbol{\sim_\mu^{\mathcal P(\Sigma^\ast)}}$ on languages over $\Sigma^\ast$ by

$L_1\sim_\mu^{\mathcal P(\Sigma^\ast)} L_2
\quad\Longleftrightarrow\quad
\mathcal R^L(L_1,L_2)\ge\mu.
$
\end{definition}
We can now introduce the appropriate quotient notions.
\begin{definition}Let
$\Sigma^\mu=\Sigma/{\sim_\mu^\Sigma}$
be the {\em quotient alphabet} with quotient map
$q_\mu:\Sigma\to\Sigma^\mu$ given by $q_\mu(a)=[a]_\mu$, and define:\\
$\bullet$ the extension of $q_\mu$ to words by putting $
q_\mu^\ast(a_1\cdots a_n)
=
[a_1]_\mu\cdots[a_n]_\mu
$; \\
$\bullet$ the quotient image of $L\subseteq\Sigma^\ast$ by putting
$ L^\mu=
\{q_\mu^\ast(w)\mid w\in L\}
\subseteq(\Sigma^\mu)^\ast.
$
\end{definition}
We now give the main definitions and results.

\begin{definition}[Fuzzy language derivative]
For $a\in\Sigma$ and $L\subseteq\Sigma^\ast$, define\\
\ \ \ \ $
\partial_a^\mu(L)
=
\{w\in\Sigma^\ast\mid
\exists b\in\Sigma,\exists v\in\Sigma^\ast:
\mathcal R(a,b)\ge\mu,\ 
\mathcal R^w(w,v)\ge\mu,\ 
bv\in L
\}.
$
\end{definition}

\begin{definition}[Fuzzy expression derivative]
For $a\in\Sigma$ and a regular expression $r$, define
$\mathcal D_a^\mu(r)=\sum\limits_{\substack{a\sim_\mu^\Sigma b\\ s\sim_\mu^e r}}\mathcal D_b(s)$,
where $\mathcal D_b(s)$ is the ordinary Brzozowski derivative.
\end{definition}
\begin{theorem}[Regularity preservation] For every symbol $a\in\Sigma$ and $0<\mu\le1$, if $L\subseteq\Sigma^\ast$ is regular, then
$\partial_a^\mu(L)$ is regular, and if $r$ is a regular expression, then $\mathcal D_a^\mu(r)$ is a regular expression.
\end{theorem}

\begin{theorem}[Finiteness]\label{finiteness}
If $\Sigma$ is finite and $L\subseteq\Sigma^\ast$ is regular, then for every word $u\in \Sigma^\ast$ and $0<\mu\le1$, the set
$
\{\partial_u^\mu(L)\mid u\in\Sigma^\ast\}
$
is finite, where $\partial_\varepsilon^\mu(L)=L$ and $\partial_{au}^\mu(L)=\partial_{u}^\mu(\partial_{a}^\mu(L))$.
\end{theorem}

\begin{theorem}[Semantic--syntactic correspondence]
For every regular expression $r$, symbol \mbox{$a\in\Sigma$}, and cut value
$0<\mu\le1$, $
L(\mathcal D_a^\mu(r))
=
\partial_a^\mu(L(r)).
$
\end{theorem}

\begin{proposition} For all restricted regular expressions $r,s$,\\
$\bullet$ 
$
\mathcal R^e(r,s)
\le
\mathcal R^L\bigl(L(r),L(s)\bigr).
$\\
$\bullet$ $\mathcal R^L\bigl(L(r),L(s)\bigr)
= \sup_{\substack{L(r)=L(r')\\
L(s)=L(s')}}\mathcal R^e(r',s').$
\end{proposition}

Finally, Theorem~\ref{main} allows us to show that ordinary derivative
automata, over the quotient alphabet $\Sigma^\mu$, capture fuzzy derivatives.


\begin{theorem}\label{main} If $a\in\Sigma$ and $L\subseteq\Sigma^\ast$ is regular, then for every cut value $0<\mu\le1$, we have \mbox{$
\partial_a^\mu(L)
=
\{w\in\Sigma^\ast\mid
q_\mu^\ast(w)\in
\partial_{[a]_\mu}(L^\mu)\}.
$}
\end{theorem}


Most of our results are formalized in Rocq~\cite{Rocq}, see~\cite{OurRocqRepo}. We build on the Coquand--Siles~\cite{CoquandRegexEquality,RocqRegexpBrzozowski} formalization of
regular-expression equivalence via Brzozowski derivatives. Theorem~\ref{finiteness} is crucial in that not all notions of finiteness are equivalent in a constructive setting. We are currently working on proving a Kuratowski-finiteness result by using normalized representatives of regular
expressions, so that repeated derivatives can be checked to generate only finitely many distinct
states using the equivalence procedure in~\cite{CoquandRegexEquality}. This should be a slight improvement, since the result in~\cite{CoquandRegexEquality} is based on a weaker constructive notion of finiteness.

\begin{thebibliography}{10}

\bibitem{AlvesKesnerRamos2025}
Sandra Alves, Delia Kesner, and Miguel Ramos.
\newblock Extending the quantitative pattern-matching paradigm.
\newblock In Oleg Kiselyov, editor, \emph{Programming Languages and Systems},
pages 84--105.
\newblock Springer Nature Singapore, Singapore, 2025.

\bibitem{B64}
Janusz A. Brzozowski.
\newblock Derivatives of regular expressions.
\newblock \emph{Journal of the ACM}, 11(4):481--494, 1964.

\bibitem{C71}
John H. Conway.
\newblock \emph{Regular Algebra and Finite Machines}.
\newblock Chapman and Hall, London, 1971.

\bibitem{CoquandRegexEquality}
Thierry Coquand and Vincent Siles.
\newblock A decision procedure for regular expression equivalence in type theory.
\newblock In \emph{Proceedings of the 1st International Conference on Certified Programs and Proofs (CPP)}, 2011.

\bibitem{DHMM}
Besik Dundua, Furio Honsell, Mariam Katamashvili, and David Maisuradze.
\newblock \emph{Fuzzy Automata}.
\newblock In preparation.

\bibitem{OurRocqRepo}
Mariam Katamashvili and David Maisuradze.
\newblock Rocq formalization for regular-expression derivatives with similarity.
\newblock \texttt{https://github.com/DUT1K1/regular-expression-derivative}.

\bibitem{met05}
George Metcalfe.
\newblock Fundamentals of fuzzy logics.
\newblock Lecture notes, University of Technology, Vienna, 2005.
\newblock \texttt{https://www.logic.at/tbilisi05/Metcalfe-notes.pdf}.

\bibitem{OwensReppyTuron}
Scott Owens, John Reppy, and Aaron Turon.
\newblock Regular-expression derivatives re-examined.
\newblock \emph{Journal of Functional Programming}, 19(2):173--190, 2009.

\bibitem{RocqRegexpBrzozowski}
Rocq community.
\newblock \emph{regexp-Brzozowski}.
\newblock \texttt{https://github.com/rocq-community/regexp-Brzozowski}.

\bibitem{Rocq}
The Rocq Prover Development Team.
\newblock The Rocq Prover, version 9.1.0.
\newblock \texttt{https://rocq-prover.org/releases/9.1.0}, 2025.

\bibitem{Sessa2002}
Maria I. Sessa.
\newblock Approximate reasoning by similarity-based SLD resolution.
\newblock \emph{Theoretical Computer Science}, 275(1--2):389--426, 2002.

\bibitem{SilvaBonchiBonsangueRutten2011}
Alexandra Silva, Filippo Bonchi, Marcello Bonsangue, and Jan Rutten.
\newblock Quantitative Kleene coalgebras.
\newblock \emph{Information and Computation}, 209(5):822--849, 2011.
\newblock Special issue: 20th International Conference on Concurrency Theory
(CONCUR 2009).

\end{thebibliography}

\end{document}

